\section{Waiting for the first success}

All examples so far have used finite sample spaces. We now return to the coin
model, but change the experiment. Instead of deciding in advance how many tosses
will be made, we toss until a certain event happens. This produces a sample
space that is infinite but still discrete.

This example is important because it shows that a sample space does not have to
be finite. It also shows why a probability model must describe not only the
possible outcomes, but also the rule that generates them.

\subsection*{The experiment}

Suppose a coin is tossed repeatedly until the first head appears. We stop as
soon as the first head occurs. Assume
\[
P(H)=p,\qquad P(T)=1-p,\qquad 0<p\leq 1.
\]
Here heads will be interpreted as success and tails as failure.

The physical experiment is not a fixed number of tosses. It may stop after one
toss, or after two tosses, or after three tosses, and so on. There is no largest
possible stopping time.

\subsection*{Recording the whole sequence}

One natural recording rule is to record the whole sequence observed up to and
including the first head. Then the possible elementary outcomes are
\[
H,\quad TH,\quad TTH,\quad TTTH,\quad \ldots
\]
The sample space is
\[
\Omega=\{H,TH,TTH,TTTH,\ldots\}.
\]
Equivalently,
\[
\Omega=\{T^{k-1}H:k=1,2,3,\ldots\},
\]
where \(T^{k-1}H\) means \(k-1\) tails followed by one head.

For example:
\[
k=1 \quad\Rightarrow\quad H,
\]
\[
k=2 \quad\Rightarrow\quad TH,
\]
\[
k=5 \quad\Rightarrow\quad TTTTH.
\]

This sample space is countably infinite. Its elements can be listed in a
sequence, but the list never ends.

\subsection*{Recording the stopping time}

Another natural recording rule is to record only the number of the toss on which
the first head appears. Let
\[
N=\text{the number of tosses needed to obtain the first head}.
\]
Then the sample space is
\[
\Omega' = \{1,2,3,\ldots\}.
\]
The outcome \(N=1\) corresponds to the sequence \(H\). The outcome \(N=2\)
corresponds to \(TH\). The outcome \(N=k\) corresponds to
\[
\underbrace{TT\cdots T}_{k-1\text{ tails}}H.
\]

The two descriptions are equivalent for this experiment. Recording the full
stopped sequence and recording the stopping time contain the same information,
because once the stopping time is known, the sequence must consist of failures
before the last toss and a success at the last toss.

\begin{remarkbox}
The sample space \(\{1,2,3,
\ldots\}\) is not the set of all possible numbers in the world. It is the set of
possible stopping times for this particular experiment.
\end{remarkbox}

\subsection*{Deriving the probability formula}

We now derive the probability that the first head appears on toss \(k\). The
condition
\[
N=k
\]
means that the first \(k-1\) tosses are tails and the \(k\)-th toss is heads.
Thus the corresponding sequence is
\[
\underbrace{TT\cdots T}_{k-1\text{ times}}H.
\]

The probability of one tail is \(1-p\), and the probability of one head is
\(p\). Therefore
\[
P(N=k)=\underbrace{(1-p)(1-p)\cdots(1-p)}_{k-1\text{ factors}}p.
\]
Hence
\[
P(N=k)=(1-p)^{k-1}p,\qquad k=1,2,3,\ldots.
\]

\begin{theorembox}
If independent tosses are repeated until the first success, and each toss has
success probability \(p\), then the probability that the first success occurs on
trial \(k\) is
\[
P(N=k)=(1-p)^{k-1}p,\qquad k=1,2,3,\ldots.
\]
\end{theorembox}

This is called the geometric model. In this chapter we do not need the name as
much as the construction. The formula appears because the outcome \(N=k\)
requires exactly \(k-1\) failures followed by one success.

\subsection*{Examples}

\begin{examplebox}
Suppose the coin is fair, so \(p=1/2\). Then
\[
P(N=1)=\frac12,
\]
\[
P(N=2)=\frac12\cdot\frac12=\frac14,
\]
\[
P(N=3)=\frac12\cdot\frac12\cdot\frac12=\frac18,
\]
and in general
\[
P(N=k)=\left(\frac12\right)^k.
\]
For a fair coin, the probability halves each time we delay the first head by one
more toss.
\end{examplebox}

\begin{examplebox}
Suppose \(P(H)=0.3\). Then \(P(T)=0.7\), and
\[
P(N=4)=0.7^3\cdot0.3.
\]
This is the probability of the sequence
\[
TTTH.
\]
The first three tosses must fail, and the fourth toss must succeed.
\end{examplebox}

\subsection*{Checking that the probabilities add to one}

A probability assignment on a countable sample space should assign total
probability one to the whole space. Here the elementary probabilities are
\[
p,\quad (1-p)p,\quad (1-p)^2p,\quad (1-p)^3p,\quad \ldots
\]
Their sum is
\[
p+(1-p)p+(1-p)^2p+(1-p)^3p+\cdots.
\]
Factor out \(p\):
\[
p\left[1+(1-p)+(1-p)^2+(1-p)^3+\cdots\right].
\]
The expression in brackets is a geometric series with ratio \(1-p\). Since
\(0<p\leq 1\), we have \(0\leq 1-p<1\), and
\[
1+(1-p)+(1-p)^2+\cdots=\frac1{1-(1-p)}=\frac1p.
\]
Therefore
\[
p\cdot\frac1p=1.
\]
Thus the probabilities of all possible stopping times add to one.

\begin{remarkbox}
This calculation also explains why we require \(p>0\). If \(p=0\), success never
occurs, and the experiment ``wait until the first success'' does not produce a
finite stopping time.
\end{remarkbox}

\subsection*{Events involving the stopping time}

Once the stopping time model is built, events are subsets of
\[
\{1,2,3,\ldots\}.
\]
For example, the event that the first success occurs no later than the third
trial is
\[
\{N\leq 3\}=\{1,2,3\}.
\]
Its probability is
\[
P(N\leq 3)=p+(1-p)p+(1-p)^2p.
\]
There is also a faster way to compute it. The complement of \(\{N\leq 3\}\) is
\[
\{N>3\},
\]
which means that the first three trials are all failures. Hence
\[
P(N>3)=(1-p)^3.
\]
Therefore
\[
P(N\leq 3)=1-(1-p)^3.
\]

More generally,
\[
P(N\leq m)=1-(1-p)^m.
\]
This is the probability that at least one success occurs in the first \(m\)
trials.

Similarly,
\[
P(N>m)=(1-p)^m,
\]
because the first \(m\) trials must all be failures.

\subsection*{Waiting for the \(r\)-th success}

A natural extension is to wait not for the first success, but for the \(r\)-th
success. Let
\[
N_r=\text{the trial on which the }r\text{-th success occurs}.
\]
For \(N_r=n\) to happen, the \(n\)-th trial must be a success, and among the
first \(n-1\) trials there must be exactly \(r-1\) successes.

The first \(n-1\) trials may contain those \(r-1\) successes in
\[
\binom{n-1}{r-1}
\]
ways. Each full sequence with \(r\) successes and \(n-r\) failures has
probability
\[
p^r(1-p)^{n-r}.
\]
Therefore
\[
P(N_r=n)=\binom{n-1}{r-1}p^r(1-p)^{n-r},
\qquad n=r,r+1,r+2,\ldots.
\]

\begin{examplebox}
Suppose we wait for the second head. The event \(N_2=5\) means that the fifth
toss is heads and exactly one of the first four tosses is heads. The position of
that earlier head can be chosen in
\[
\binom41=4
\]
ways. Hence
\[
P(N_2=5)=\binom41 p^2(1-p)^3.
\]
\end{examplebox}

This extension is included to show that the first-success model is not an
isolated trick. It is the first member of a family of waiting-time models. The
logic is always the same: describe what must happen before the stopping trial,
describe what must happen at the stopping trial, count the possible positions of
earlier successes, and multiply the appropriate probabilities.

\subsection*{What waiting-time models teach}

The waiting-time model teaches several new lessons.

First, a sample space may be infinite. The set \(\{1,2,3,\ldots\}\) is a valid
discrete sample space when it records a stopping time.

Second, an infinite sample space still requires probabilities that add to one.
In the first-success model, this is guaranteed by the geometric series.

Third, the elementary outcome may be represented in different but equivalent
ways. We may record the stopped sequence \(T^{k-1}H\), or we may record the
number \(k\). The choice depends on which description is more useful.

Fourth, a formula such as
\[
P(N=k)=(1-p)^{k-1}p
\]
should be read as a description of the required pattern: failures before the
stopping trial and success at the stopping trial.

\begin{remarkbox}
Waiting-time models are a bridge between elementary finite probability and later
probability distributions. They show that probability spaces can be built from a
process that stops at a random time, and that even an infinite model can be
handled by careful construction from elementary outcomes.
\end{remarkbox}
